%! Unit 1, Lessons 1.0–1.4 — SAMPLE Quiz (blank)
% The practice form of the mid-unit quiz. Same shape, same item count, same point spread as
% ../actual_quiz/ — a different context in every item, so working this one teaches the moves
% rather than the answers. Handed out with ../study_guide/ and worked before the real quiz.
% NAMESTRIP: no \namedateperiod — this one is practice taken home, not an assessment sat in a
% testing setting. The ACTUAL quiz carries the name row; this one must not.
% PAGE LOCKSTEP with ../../quiz_keys/sample_quiz_key/: every \writelines{n} here is answered by
% exactly n \ansline{} there, each terse enough not to wrap. Nothing else differs between the
% two files except the preamble, the header, and the \ans{} marks on the correct MC options.
\documentclass[12pt]{article}
\usepackage{apstats-article}
\usepackage{apstats-boxes}

\begin{document}

\pageheader{Unit 1 \quad Lessons 1.0--1.4}{Sample Quiz}

\vspace{0.14in}

\boxguard[8]
\begin{remindbox}
\small
This is the \textbf{practice} quiz --- it is not scored. Work it once with your notes closed,
then check it against the key. It has the same number of items, in the same order, as the quiz
you will sit. Every free-response answer must name the variable and its units.
\end{remindbox}

\vspace{0.12in}

\boxguard[24]
\begin{scenariobox}[The Riverside Farmers Market]{navy}
\small
The Riverside Farmers Market keeps one row of records for each of its \textbf{48 vendors}. Four
of the 48 rows are shown below.

\vspace{6pt}
\begin{center}
\renewcommand{\arraystretch}{1.2}
\begin{tabular}{@{} l l c r @{}}
\rowcolor{sky}
\textbf{Vendor} & \textbf{Stall type} & \textbf{Stall number} & \textbf{Weekly revenue} \\
\midrule
Okonkwo Gardens   & Produce        & 104 & \$620 \\
Hollis Bakehouse  & Baked goods    & 117 & \$480 \\
Pine Street Clay  & Crafts         & 129 & \$310 \\
Sopa y Sal        & Prepared food  & 141 & \$755 \\
\end{tabular}
\end{center}

\vspace{6pt}
Across all 48 vendors, the stall types are: \textbf{Produce 18}, \textbf{Baked goods 12},
\textbf{Crafts 10}, \textbf{Prepared food 8}.
\end{scenariobox}

\vspace{0.14in}

\boxguard[16]
\begin{headlinebox}{goldbg}
\small
\textbf{Section I --- Multiple Choice.} Circle the single best answer. Read all five options ---
more than one of them is about the context, not just the arithmetic.
\end{headlinebox}

\vspace{0.10in}

\begin{enumerate}[leftmargin=1.9em, itemsep=0.13in]

  \item In the market's records, the \textbf{individuals} are
        \begin{enumerate}[label=\textbf{(\Alph*)}, leftmargin=2.2em, itemsep=2pt]
          \item the four stall types.
          \item the 48 vendors.
          \item the weekly revenues.
          \item the market itself.
          \item the shoppers who visit the market.
        \end{enumerate}

  \item Which of the following is a \textbf{statistical question} about these vendors?
        \begin{enumerate}[label=\textbf{(\Alph*)}, leftmargin=2.2em, itemsep=2pt]
          \item How much did Hollis Bakehouse earn last week?
          \item How many stalls does the market have?
          \item What is stall 129's stall type?
          \item How much does weekly revenue vary among the market's 48 vendors?
          \item Is the market open on Saturdays?
        \end{enumerate}

  \item \textbf{Stall number} is best classified as
        \begin{enumerate}[label=\textbf{(\Alph*)}, leftmargin=2.2em, itemsep=2pt]
          \item quantitative, because its values are written in digits.
          \item quantitative and discrete, because the numbers are whole numbers.
          \item quantitative and continuous, because 117.5 is a number.
          \item categorical, because the number labels a stall rather than recording an amount.
          \item categorical, because there are only 48 of them.
        \end{enumerate}

  \item Bayview Market has 30 vendors, 15 of them selling produce. Which statement about the
        two markets is correct?
        \begin{enumerate}[label=\textbf{(\Alph*)}, leftmargin=2.2em, itemsep=2pt]
          \item Riverside is more produce-focused, because 18 is greater than 15.
          \item Bayview is more produce-focused, because 50\% is greater than 37.5\%.
          \item The two markets are equally produce-focused.
          \item No comparison is possible, because the markets have different numbers of vendors.
          \item Bayview is more produce-focused, because it has fewer vendors overall.
        \end{enumerate}

  \item \textbf{Pounds of tomatoes sold in a week}, recorded to the nearest whole pound, is
        \begin{enumerate}[label=\textbf{(\Alph*)}, leftmargin=2.2em, itemsep=2pt]
          \item categorical, because the amounts were rounded.
          \item discrete, because every recorded value is a whole number.
          \item continuous, because between any two weights another weight is possible.
          \item discrete, because you can count the tomatoes.
          \item neither, because weight cannot be graphed.
        \end{enumerate}

  \item A dotplot of weekly revenue piles up at the low end and trails off toward the high end,
        with a few vendors far to the right. This distribution is
        \begin{enumerate}[label=\textbf{(\Alph*)}, leftmargin=2.2em, itemsep=2pt]
          \item skewed left, because most of the dots sit on the left.
          \item skewed right, because the longer tail runs out to the right.
          \item symmetric, because it has one peak.
          \item uniform, because every vendor earns something.
          \item impossible to name without the center.
        \end{enumerate}

\end{enumerate}

\vspace{0.14in}

% No \newpage here: the guard keeps the Section II strip off a page edge, and letting
% the section flow saves the packet a near-empty final page.
\boxguard[10]
\begin{headlinebox}{goldbg}
\small
\textbf{Section II --- Free Response.} Show your reasoning. Every answer must be written
\emph{in the context of the problem} --- name the variable and its units.
\end{headlinebox}

\vspace{0.10in}

\begin{enumerate}[leftmargin=1.9em, itemsep=0.16in, resume]

  \item Use the Riverside Farmers Market records above.

        \vspace{6pt}
        \textbf{(a)} Identify the individuals, then name one categorical variable and one
        quantitative variable in the table.\\[4pt]
        \writelines{2}

        \vspace{6pt}
        \textbf{(b)} Give the relative frequency of each of the four stall types, to the nearest
        tenth of a percent.\\[4pt]
        \writelines{1}

        \vspace{6pt}
        \textbf{(c)} A reporter writes, ``Riverside is the region's produce market --- it has
        more produce vendors than Bayview does.'' Riverside has 18 of 48; Bayview has 15 of 30.
        Explain what is wrong with the reporter's reasoning, and state which market a share
        comparison favours.\\[4pt]
        \writelines{3}

        \vspace{6pt}
        \textbf{(d)} Write one statistical question that could be answered using these records,
        and say which of your two checks makes it statistical.\\[4pt]
        \writelines{2}

  \item The dotplot below shows the \textbf{weekly revenue}, in dollars, of all 30 vendors at a
        neighbouring market last week.

        \vspace{4pt}
        \noindent\hfil
        \begin{tikzpicture}[scale=0.75]
          \draw[->, thick] (-0.2,0) -- (10.6,0);
          \foreach \x/\lab in {0/0, 2/400, 4/800, 6/1200, 8/1600, 10/2000} {
            \draw (\x,-0.07) -- (\x,0.07); \node[below, font=\tiny] at (\x,-0.1) {\lab};}
          \foreach \x in {1, 3, 5, 7, 9} { \draw (\x,-0.05) -- (\x,0.05); }
          \foreach \x/\n in {1/2, 1.5/5, 2/6, 2.5/5, 3/4, 3.5/3, 4/2, 4.5/1, 5/1, 9.5/1} {
            \foreach \k in {1,...,\n} { \fill[navy] (\x,{0.2*\k}) circle (0.07); }}
          \node[font=\tiny] at (5.2,-0.85) {Weekly revenue (dollars)};
        \end{tikzpicture}\hfil
        \vspace{6pt}

        \textbf{(a)} Is weekly revenue discrete or continuous? Justify from the variable itself,
        not from the display.\\[4pt]
        \writelines{2}

        \vspace{6pt}
        \textbf{(b)} Describe the distribution of weekly revenue completely.\\[4pt]
        \writelines{4}

        \vspace{6pt}
        \textbf{(c)} A classmate writes, ``It is skewed left, because most of the dots are on the
        left side.'' Correct the classmate, and say what actually decides the name.\\[4pt]
        \writelines{2}

        \vspace{6pt}
        \textbf{(d)} The market's manager proposes dropping the vendor near \$1{,}900 before
        reporting a typical revenue, calling it ``clearly a mistake.'' Give one reason that is a
        poor decision, and say what should be done with that vendor instead.\\[4pt]
        \writelines{2}

\end{enumerate}

\end{document}
